\section{Discussion and Concluding Remarks}

We have demonstrated that the long-standing trade-off between \textit{semantic abstraction} and \textit{spatial grounding} (the Invariance Paradox) is largely an artifact of the traditional dense-grid representation. By factorizing the latent space into sparse ``What'' and ``Where'' components, STELLAR effectively resolves this conflict.

The core of STELLAR’s success lies in the principle of \textit{semantic triage}. In a dense model (e.g., MAE), the representation is spatially exhaustive but semantically diluted; the model is forced to allocate representational capacity to every patch, including stochastically redundant background noise. 


The low-rank factorization provides the mathematical machinery to disentangle two distinct types of visual information. The \textit{Concept Tokens} ($\mathbf{S}$) are trained to be view-invariant, capturing the categorical ``What,'' while the \textit{Spatial Coefficients} ($\mathbf{L}$) remain equivariant, capturing the geometric ``Where.'' 

This disentanglement allows the model to satisfy the Joint Embedding objective (alignment across views) without destroying the spatial anchors needed for high-fidelity reconstruction. By separating these factors, we avoid the ``semantic blurring'' often seen in global-pooling methods, as each sparse token maintains a precise, albeit flexible, relationship with the physical image geometry.

\textbf{The Path to Unified Multimodality}: Perhaps the most promising frontier for STELLAR is its potential as a visual front-end for Large Language Models (LLMs). Because our tokens are sparse and semantically grounded, they offer a more natural interface for cross-modal alignment than the hundreds of dense tokens generated by standard ViTs. Future work will investigate the systematic integration of STELLAR concepts with linguistic embeddings for more interpretable multimodal understanding.

Overall, our results highlight sparse tokens as a promising direction for unifying efficiency, interpretability, and semantic richness in self-supervised representation learning.

% \section{Discussion and Conclusion}

% We demonstrate that sparse visual representation, when structured in form of low-rank factorization, can simultaneously support high-fidelity reconstruction and rich semantic understanding. By disentangling \textit{what} and \textit{where}, STELLAR learns compact tokens that transfer effectively to both global and fine-grained tasks, surpassing prior methods. The sparse image modeling framework effectively shapes the feature map into semantic regions, without direct loss on the dense representation.

% Nevertheless, several limitations remain. First, our experiments are conducted primarily on ImageNet-1K scale; extending to larger pretraining corpora is likely to further narrowing the gap with large-scale models such as DINOv2. Second, the current framework adopts a minimalist architecture. Exploring more expressive decoders or hybrid architectures could enhance reconstruction fidelity. Finally, while STELLAR naturally provides a compact interface for multimodal alignment, we only experimented simple alignment in appendix, leaving systematic integration with language models to future work. Overall, our results highlight sparse tokens as a promising direction for unifying efficiency, interpretability, and semantic richness in representation learning.


\section*{Impact Statement}

This paper presents work whose goal is to advance the field of machine learning. There are many potential societal consequences of our work, none of which we feel must be specifically highlighted here.